\section{Lecture — 24/09/2026}


\subsection*{Curriculum Position}
\begin{itemize}
	\item Week-1 — Course intro, syllabus, Newton's laws, vectors.
\end{itemize}


\subsection*{Key Content}
\begin{enumerate}
	\item Mechanics 
	\item Basic Quantities 
	\item Idealizations and an intro into rigid body mechanics 
	\item Newton's Laws of Motion
	\item Newton's Law of Gravitational Attraction 
	\item SI Units
	\item Unit Prefixes and Rules for Unit Prefixes
	\item Dimensional Homogeneity
	\item Sig-Figs and Rounding
	\item General Procedure for Analyzing Systems
\end{enumerate}


\subsection*{Instructor Emphasis}
\begin{itemize}
	\item Use 3 Sig-Figs for calculations.
	\item Take g = 9.81 $ \text{m/s}^2 $ for calculations.
	\item He cares a lot about unit prefix rules.
	\item It is recommended that calculations be done with the units alongside, \\
	so that the resulting unit gets verified.
	\item Retain a greater number of digits than the problem data for calculations.
	\item For solving problems with different units, converting to base units \\ 
	first is recommended 
\end{itemize}


\subsection*{Announcements}
\begin{enumerate}
	\item Email appointment before office visit.
	\item Lecture notes are provided and said to be adequate.
	\item Nonattendance allowed up to 4 days.
	\item AA is (100-90) by default, Grading curve is possible based on performance.
	\item Worked out problems will be provided on the course Teams group.
	\item 1 Midterm (35\%), 1 Final (45\%), Pop-Quizzes (5\% each) 
	Announced a week before.  
	\item Textbook: Hibbeler, R.C. Engineering Mechanics
	\item Next week will consist of 3 sessions.
\end{enumerate}


\subsection*{Open Questions}

\subsection*{Independent Follow-up}
\begin{itemize}
	\item Minimum: Have a light read.
	\item Exploration: Watch some stuff on statics and FEA
\end{itemize}
